\section{Ordered and unordered opponents}
Assume common opponent action labels, and write $h(t)$ for an opponent tuple's
histogram. Ex post indifference to opponent arrangements means
$u_i(a,t)=w_i(a,h(t))$. This holds exactly when every opponent permutation
fixing the focal player preserves that focal payoff: two tuples have the same
histogram exactly when a permutation carries one to the other. This is a
factorization statement, not a consequence of concealing identities.

If instead beliefs are uniform conditional on the histogram, the evaluated
table is the average of $u_i$ on each histogram orbit. Equality under every
exchangeable opponent law detects only these averages: uniform laws on each
orbit prove necessity, and decomposition into orbit-uniform laws proves
sufficiency. Expected averaging requires a selected cardinal representation.

For cyclic indices on three binary players, take
$u_i(s)=s_i+s_{i+1}-s_{i+2}$. Conditional on a histogram, each focal average is
its own action $s_i$. Yet opponent arrangements $(1,0)$ and $(0,1)$ give
payoffs $1$ and $-1$ when the focal action is 0. Aggregate tests therefore
miss genuine ordered-opponent effects. Allowing strategy relabelings changes
the admissible physical comparisons and must be specified separately.
